% ============================================================
% GEMEINSAME ARCHITEKTUR-SECTION
% Empfehlung: In deiner architecture.tex oder als eigenes
% Kapitel vor den technischen Einzelteilen
% ============================================================

\section{Gesamtarchitektur der Anwendung}
\setauthor{Markus Tran}

Die vorliegende Anwendung folgt einer klar getrennten Client-Server-Architektur,
bei der Frontend und Backend als eigenständige Einheiten entwickelt wurden und
ausschließlich über definierte REST-Schnittstellen miteinander kommunizieren.

Das Frontend wurde von \textit{Abdulkerim Cufurovic} als Single-Page-Application
mit Angular umgesetzt und übernimmt die gesamte Darstellung, Navigation und
Benutzerinteraktion. Das Backend, dessen Implementierung Gegenstand des vorliegenden
Teils der Diplomarbeit ist, wurde mit Quarkus realisiert und stellt die
Geschäftslogik, Datenbankanbindung sowie die Kommunikation mit der externen
TMDB API bereit.

Die Authentifizierung beider Schichten erfolgt über Keycloak als zentralen
Identity-Provider. Dieses Zusammenspiel ist in Abbildung~\ref{fig:gesamtarchitektur}
dargestellt.

\begin{figure}[h]
    \centering
    \includegraphics[width=0.95\textwidth]{images/architektur_gesamt}
    \caption{Gesamtarchitektur der Anwendung mit Frontend, Backend und Keycloak}
    \label{fig:gesamtarchitektur}
\end{figure}

Die strikte Trennung zwischen Frontend und Backend bietet mehrere Vorteile.
Beide Schichten können unabhängig voneinander weiterentwickelt, getestet und
deployed werden. Änderungen an der Benutzeroberfläche erfordern keine Anpassungen
im Backend und umgekehrt, solange die vereinbarten Schnittstellen eingehalten
werden. Diese Architekturentscheidung war für die Aufteilung der Diplomarbeit
in zwei eigenständige Projektteile grundlegend.